\section*{Erd\H{o}s Problem 453}

\paragraph{Statement and result.}
Must every sufficiently large \(n\) admit \(1\leq i<n\) such that
\(p_n^2<p_{n-i}p_{n+i}\)? No. We reconstruct Pomerance's supporting-line
argument, obtaining infinitely many \(n\) for which all these inequalities
are reversed strictly.

\paragraph{A general sequence lemma.}
Let \(a_n\) be an increasing unbounded sequence with \(a_n=o(n)\).
For \(t>0\), the sequence \(a_j-tj\) tends to \(-\infty\), so it has
a maximum. Exclude the countable set of slopes
\[
 \left\{\frac{a_j-a_k}{j-k}:j\ne k\right\}.
\]
For every remaining \(t\), the maximizing index \(n(t)\) is unique.
Consequently, for every \(1\leq i<n(t)\),
\[
 a_{n(t)-i}-t(n(t)-i)<a_{n(t)}-tn(t),\qquad
 a_{n(t)+i}-t(n(t)+i)<a_{n(t)}-tn(t).
\]
Adding cancels the linear terms and gives
\[
 a_{n(t)-i}+a_{n(t)+i}<2a_{n(t)}.
 \tag{1}
\]

These maximizing indices are unbounded as \(t\downarrow0\) through
allowed slopes. Indeed, for any \(M\), choose \(K>M\) with
\(a_K>\max_{j\leq M}a_j\). For sufficiently small positive \(t\),
\(a_K-tK>\max_{j\leq M}(a_j-tj)\). Thus \(n(t)>M\).
There are allowed slopes arbitrarily close to zero, since a countable
set cannot exhaust any real interval. This proves that infinitely many
indices satisfy (1).

\paragraph{Application to primes.}
Take \(a_n=\log p_n\). Infinitude of primes makes this sequence
increasing and unbounded. The standard Chebyshev bound
\(p_n=O(n\log n)\) gives \(\log p_n=O(\log n)=o(n)\).
For the infinitely many indices supplied by the lemma, exponentiating
(1) yields
\[
 p_{n-i}p_{n+i}<p_n^2\qquad(1\leq i<n).
\]
This disproves the proposed eventual assertion.

\paragraph{Attribution and scope.}
The supporting-line argument is a reconstruction of the proof attributed
to C.~Pomerance (1979) at \url{https://www.erdosproblems.com/453},
checked 24 September 2026. Chebyshev's prime bound is the only
number-theoretic external input. No claim about the density of these
indices is needed or made.

\paragraph{Completion estimate.}
The full negative answer, including simultaneous strict inequalities, is proved.
\textbf{COMPLETION: 100\%}
